% sec-03: three instruments on eight attributes.  Figure 4 (d2-type grid),
% Table 8.  No amount appears in the figure, by the rule stated in sec-00.
\section{Three Instruments, One Deciding Row}

Three instruments could raise the same amount. They differ on eight attributes,
and seven of those eight are trade-offs a founder can weigh from experience. The
eighth is not a trade-off at all: it is a federal eligibility question with a
right answer, and it decides the choice for a company in this position.

\begin{appfloat}[!tb]
\begin{appfig}[d2-type / three-column container comparison, no edges]
% Nine rows on a 0.68cm pitch.  The three instrument columns are equal at 26mm
% so that no column reads as favored by width; the attribute column is 40mm
% because its labels are the longest cells in the figure.  There are no edges:
% the three instruments are alternatives, not stages.
\node[d2cellh,minimum width=40mm,text width=40mm] at (0,0) {Attribute};
\foreach \x/\lab in {3.15/{Post-money SAFE}, 5.85/{Convertible note}, 8.55/{Priced preferred}}{
  \node[d2cellh,minimum width=26mm,text width=24mm] at (\x,0) {\lab};}
\foreach \y/\a/\s/\n/\p in {%
  -0.68/{Time to first close}/{Days}/{Days to weeks}/{Weeks to months},
  -1.36/{Legal cost to issuer}/{Lowest}/{Low}/{Highest},
  -2.04/{Accrues interest}/{No}/{Yes}/{No},
  -2.72/{Has a maturity date}/{No}/{Yes}/{No},
  -3.40/{Ownership at signing}/{None}/{None}/{Immediate},
  -4.08/{Cap table clarity later}/{Good}/{Good}/{Best},
  -4.76/{Governance given away}/{None}/{None}/{Board seat usual}}{
  \node[d2celll,minimum width=40mm,text width=40mm] at (0,\y) {\a};
  \node[d2cell,minimum width=26mm,text width=24mm] at (3.15,\y) {\s};
  \node[d2cell,minimum width=26mm,text width=24mm] at (5.85,\y) {\n};
  \node[d2cell,minimum width=26mm,text width=24mm] at (8.55,\y) {\p};}
\draw[fundmid,line width=0.9pt] (-1.05,-5.10) -- (9.85,-5.10);
\node[d2celll,minimum width=40mm,text width=40mm] at (0,-5.44) {SBIR ownership answer};
\node[d2cellk,minimum width=26mm,text width=24mm] at (3.15,-5.44) {Unchanged};
\node[d2cellk,minimum width=26mm,text width=24mm] at (5.85,-5.44) {Unchanged};
\node[d2cellg,minimum width=26mm,text width=24mm] at (8.55,-5.44) {Can change};
\ptitle{-1.05}{0.62}{Three alternatives, not three stages}
\node[font=\tiny\itshape,text=fundgray,anchor=north west,text width=118mm]
  at (-1.05,-5.95) {No valuation, cap, discount, minimum, or closing date appears
  in this figure. Those are offering terms, and no offering exists.};
\end{appfig}
\vspace{-0.60cm}
\figcaption{Figure 4. Three candidate instruments across eight attributes at one raise\\
size, with the row that decides the choice for this company marked in fill.}
\end{appfloat}

\subsection{Why the last row dominates}

Under 13 CFR 121.702, majority ownership by multiple venture capital operating
companies, hedge funds, or private equity firms changes Small Business Innovation
Research eligibility and, at some agencies, requires a separate authorization
\cite{ecfr2026sbirownership}. The SBA Company Registry asks that question
directly, and the answer is re-checked at every financing event rather than
annually \cite{sbirregistry2026}.

An instrument that has not converted transfers no ownership, so the registry
answer on the day after signing is the same as the day before. A priced round
that transfers majority ownership can change it. That does not make a priced
round wrong; it makes it a decision that must be taken with the federal route's
status in front of the decision-maker rather than discovered afterward at a
receipt date.

That reading is the company's own and is sent to counsel for confirmation or
correction before any instrument is signed. It is item two of the engagement
letter, and it is the single most consequential item in that engagement.

\subsection{The tranche shape, if the instrument is convertible}

\begin{apptable}
\begin{tabularx}{\textwidth}{>{\raggedright\arraybackslash}p{1.9cm} >{\raggedright\arraybackslash}p{4.2cm} >{\raggedright\arraybackslash}X}
\headrow \thead{Tranche} & \thead{What it funds} & \thead{The evidence that opens it}\\
\midrule
First & Site feasibility, regulatory pre-submission, the verification harness & An instrument selected and counsel engaged \\
Second & Site agreement, IRB submission, first-participant readiness & A written site feasibility outcome from a named institution \\
Third & Trial conduct beyond what the federal route funds & An accepted investigational new drug application \\
\end{tabularx}
\end{apptable}
\vspace{-0.60cm}
\tabcap{Table 8. Three tranches keyed to evidence rather than to a calendar,\\
with the artifact that opens each stated as a condition and not a date.}

Tranching against evidence rather than against a calendar is deliberate. A holder
who funds a tranche is funding a stated next artifact, and a tranche that does not
open because its trigger did not fire is a smaller loss than a full raise spent
against a milestone that never arrived.

\subsection{The recommendation}

Take the instrument that leaves the ownership answer unchanged until a deliberate
decision is made to change it, and take it with one valuation cap rather than
several. A stack of instruments at different caps produces a cap table that a
later institutional round has to unwind before it can price, and unwinding costs
more than the difference between the caps.

That is a recommendation to the chief executive about this company's own capital
structure. It is not advice to any other person, and it is not an offer.
